\section{Weather Augmentations}
\label{sec:weather-aug}

This section addresses \textbf{RQ2}: whether input-space weather augmentations improve
real-world generalization, which mixing ratios help or harm, and whether augmentation
can compensate for a lack of synthetic diversity. It is the least direct intervention
on the domain shift --- it approximates target-domain nuisances without using any real
data.

\subsection{Augmentations and Training Sources}
\label{subsec:aug-design}
Weather augmentations are drawn from the \texttt{albumentations} library and applied
\emph{individually} to synthesize the adverse conditions of the real test set:
\textbf{glare} (RandomSunFlare), \textbf{fog} (RandomFog), \textbf{rain}
(RandomRain), \textbf{water refraction} (WaterRefraction) and \textbf{night}
(a darkening transform). A qualitative comparison of these augmentations against real
glare/fog/rain images confirms that they approximate, but do not perfectly reproduce,
the real degradations.

Following \Cref{sec:gap}, augmentations are applied to the two most relevant training
sources: \texttt{ges} (best individual source on the real set) and \texttt{combined}
(best overall). Evaluation is per weather scenario on the LARD test set. This pairing
makes it possible to test whether augmentation benefits depend on the underlying
dataset diversity.

\subsection{Controlling Sampling Variance}
\label{subsec:aug-variance}
A first sweep exposed a methodological problem: the run-to-run mAP variance was so
large that it masked any augmentation effect. For the \texttt{combined} source, an
augmentation sweep on the rain split showed \emph{no} improvement over the
no-augmentation baseline and a wide spread across the four runs; for \texttt{ges}, the
$0\%$-ratio runs had no spread at all (the training set is fixed), confirming that the
variance originates from two independent sampling steps: (i) re-sampling the
\texttt{combined} training set and (ii) sampling which images are augmented.

To attribute variance to the augmentation alone, the protocol is tightened
(\Cref{subsec:setup-protocol}):
\begin{itemize}
    \item the \texttt{combined} training set is \emph{fixed} by sampling it four times
    and keeping the highest-performing draw;
    \item for each augmentation ratio, four independent runs are performed and the
    \emph{maximum} achieved mAP is reported, i.e.\ an estimate of the best attainable
    performance at that ratio.
\end{itemize}
The ratios evaluated are \textbf{0.15, 0.30, 0.45, 0.60, 0.75}. This choice reflects
the combinatorial blow-up of a full grid: as noted in the augmentation-influence
literature~\cite{ruter2025augmentations}, sweeping augmentations, ratios and sources
jointly yields on the order of $10^9$ configurations, so the search is deliberately
scoped and, as a bridge to the later methods, the data/augmentation mix is tuned
rather than exhaustively searched.

\subsection{Results}
\label{subsec:aug-results}
\Cref{tab:aug-bestworst} summarizes, per scenario, the fixed baselines for
\texttt{ges} and \texttt{combined} together with the best- and worst-performing
augmentation$\times$ratio found. The analysis answers the three sub-questions of RQ2.

\paragraph{Do augmentations behave as expected?}
Not straightforwardly. The augmentation that matches a scenario is \emph{not}
generally the best augmentation for that scenario --- e.g.\ on the glare split the best
result for \texttt{ges} comes from water refraction, not glare, and on several splits
the scenario-matched augmentation is among the \emph{worst}. Weather robustness thus
does not decompose cleanly into ``augment with the target weather''.

\paragraph{Which ratio helps and which harms?}
The effect is strongly ratio-dependent and non-monotonic. Two representative cases on
\texttt{ges}: under water \emph{refraction} on the rain split, raising the ratio from
$0.60$ to $0.75$ improves mAP by almost $10$ points; conversely, \emph{glare}
augmentation on the fog split yields both the best and the worst outcomes, dropping by
more than $10$ mAP points when the ratio is raised from $0.45$ to $0.75$. There is no
single universally safe ratio; overly high ratios tend to over-corrupt the source.

\paragraph{Do augmentations compensate for missing diversity?}
Only weakly. For the low-diversity \texttt{ges} source, the best augmentation can
improve a scenario substantially, but it does \emph{not} reach the performance of the
\texttt{combined} baseline. And when the diversity is already high
(\texttt{combined}), augmentation helps only very little --- the \texttt{combined}
baseline is consistently at or above the best single augmentation applied to it.
Appearance augmentation therefore does not substitute for rendering-source diversity;
at best it narrows the gap for a single-source model.

\begin{table}[th]
    \caption{Weather augmentation on \texttt{ges} and \texttt{combined}, per scenario:
    fixed baseline and best/worst augmentation$\times$ratio (max over 4 runs, reduced
    20-epoch budget). The scenario-matched augmentation is rarely the best, and the
    \texttt{combined} baseline is hard to beat.}
    \label{tab:aug-bestworst}
    \centering
    \begin{NiceTabular}{llll}
        \CodeBefore
        \rowcolors{2}{gray!10}{white}
        \Body
        \toprule
        \textbf{Scenario} & \textbf{Source} & \textbf{Baseline} & \textbf{Best / Worst augmentation (ratio)} \\
        \midrule
        Glare & ges      & \tbd & refraction 15\% / glare 75\% \\
        Glare & combined & \tbd & refraction 30\% / glare 75\% \\
        Rain  & ges      & \tbd & dark 60\% / glare 75\% \\
        Rain  & combined & \tbd & fog 15\% / fog 60\% \\
        Fog   & ges      & \tbd & glare 45\% / rain 75\% \\
        Fog   & combined & \tbd & dark 15\% / fog 75\% \\
        Night & ges      & \tbd & rain 30\% / glare 75\% \\
        Night & combined & \tbd & dark 30\% / glare 45\% \\
        \bottomrule
    \end{NiceTabular}
\end{table}
\todo[inline]{Fill baseline mAP values from the fixed-dataset runs and, if desired,
the best/worst mAP alongside each configuration.}

\plotslot{GES/combined baseline with best and worst augmentation per scenario
(min--max bands); and the per-augmentation ratio curves for the \texttt{ges} rain and
fog splits (refraction $+10$ mAP at $0.75$; glare drop $>10$ mAP $0.45\!\to\!0.75$).}

\subsection{Discussion}
\label{subsec:aug-discussion}
Weather augmentation is a cheap, label-preserving baseline that requires no real data,
but its benefit here is limited and unstable: the scenario-matched augmentation is
often not the best, the useful ratio is condition-specific, and augmentation does not
recover the advantage of source diversity. Part of this instability is
methodological --- four runs are not always enough to average out the sampling
variance, and more runs are expensive --- which itself is a finding: robust
augmentation studies on this task need either many runs or a fixed, controlled data
mix. These limitations motivate the data-driven techniques of
\Cref{sec:pseudo-labels} (which uses real, unlabeled data) and \Cref{sec:style-transfer}
(which imports real appearance generatively).
